\begin{theorem}[Additivity of the index on an MPU representation]
    \label{thm:mpug_group_index_additivity}
    \lean{MPOTensor.GroupFamily.IsRawRepresentation.index_mul,
        MPOTensor.GroupFamily.IsRawRepresentation.index_one,
        MPOTensor.GroupFamily.IsRawRepresentation.neZero_phys}
    \leanok
    \uses{def:mpug_group_representation, def:mpu_reduced_representative_index}
    Let $\{\mathcal U_g\}_{g\in G}$ be an injective family of MPUs with positive
    bond dimensions, satisfying~(\ref{eq:mpug_rep_identity}) and
    (\ref{eq:mpug_rep_multiplication}) at every positive chain length.
    Simplicity is not required. Then $d>0$ and
    \begin{align}
        \operatorname{ind}(\mathcal U_e) & =0,                                                                 \\
        \operatorname{ind}(\mathcal U_{gh})
                                         & =\operatorname{ind}(\mathcal U_g)+\operatorname{ind}(\mathcal U_h).
    \end{align}
    These are the index properties used in the finite-order argument of
    \cite[line~1547]{FrancoRubio2025MPUGauging}.
\end{theorem}

\begin{proof}\leanok
    \uses{thm:mpu_reduced_representative_index_invariance,
        thm:mpu_representative_index_tensor_product,
        thm:mpu_identity_index, thm:mpu_admissible_index_composition}
    Injectivity on a positive-dimensional bond space excludes an empty
    physical alphabet, so $d>0$. The identity law identifies the periodic
    operators of $\mathcal U_e$ with those of the identity tensor, whose
    index is zero. The multiplication law identifies those of
    $\mathcal U_{gh}$ with the composition of $\mathcal U_g$ and
    $\mathcal U_h$. Invariance under equality of periodic operators and
    additivity under composition give the two conclusions.
\end{proof}

\begin{theorem}[Vanishing of the index for a finite-group representation]
    \label{thm:mpug_finite_group_index_zero}
    \lean{MPOTensor.GroupFamily.IsRawRepresentation.index_eq_zero}
    \leanok
    \uses{thm:mpug_group_index_additivity}
    Under the hypotheses of Theorem~\ref{thm:mpug_group_index_additivity},
    suppose that $G$ is finite. Then
    $\operatorname{ind}(\mathcal U_g)=0$ for every $g\in G$.
    This is the finite-order conclusion in
    \cite[line~1547]{FrancoRubio2025MPUGauging}.
\end{theorem}

\begin{proof}\leanok
    \uses{thm:mpug_group_index_additivity}
    Induction gives
    $\operatorname{ind}(\mathcal U_{g^n})=n\operatorname{ind}(\mathcal U_g)$.
    For $n=\operatorname{ord}(g)>0$, the identity $g^n=e$ gives
    $0=n\operatorname{ind}(\mathcal U_g)$; hence
    $\operatorname{ind}(\mathcal U_g)=0$.
\end{proof}

\begin{theorem}[Source-leg dimensions in a finite-group representation]
    \label{thm:mpug_finite_group_source_ranks}
    \lean{MPOTensor.GroupFamily.IsRepresentation.sourceRanks_eq_physDim}
    \leanok
    \uses{def:mpug_group_representation, def:mpu_right_rank, def:mpu_left_rank}
    Let $G$ be finite and let $\{\mathcal U_g\}_{g\in G}$ be an exact
    representation by simple injective MPU tensors of physical dimension
    $d$. For every $g\in G$, the two source-leg dimensions are
    \begin{align}
        d_r(g) & =d, \\
        d_l(g) & =d.
    \end{align}
    This is the conclusion of
    \cite[line~1547]{FrancoRubio2025MPUGauging}.
\end{theorem}

\begin{proof}\leanok
    \uses{thm:mpug_finite_group_index_zero,
        thm:mpu_simple_injective_source_index,
        thm:mpug_zero_specified_source_index_ranks}
    Put $r=d_r(g)$ and $\ell=d_l(g)$. Simplicity and injectivity give
    $r\ell=d^2$ and
    $\operatorname{ind}(\mathcal U_g)=\frac12(\log_2 r-\log_2\ell)$.
    The physical dimension is positive, so both ranks are positive.
    Finite order gives $\operatorname{ind}(\mathcal U_g)=0$. Using
    $r\ell=d^2$, the index can also be written as $\log_2(r/d)$.
    Thus $r/d=1$ and $r=d$; substitution into $r\ell=d^2$ gives $\ell=d$.
\end{proof}
